\documentclass[12pt]{article}

%% Basic document formatting
\usepackage{amsmath, amsthm, amssymb, amsfonts}
\usepackage{mathtools}
\usepackage{xspace}
\usepackage{thmtools}
\usepackage{graphicx}
\usepackage{setspace}
\usepackage{fancyhdr}
\usepackage{titling}
\usepackage{hyperref}
\usepackage[left=0.4in,right=0.4in,top=1in,bottom=1in]{geometry}
\usepackage{float}
\usepackage{tabularx}
\usepackage[utf8]{inputenc}
\usepackage[english]{babel}
\usepackage{framed}
\usepackage[dvipsnames]{xcolor}
\usepackage{environ}
\usepackage{tcolorbox}
\tcbuselibrary{theorems,skins,breakable}

\usepackage{enumitem}
\setlist[enumerate]{leftmargin=*}
\setlist[enumerate,1]{labelindent=\parindent}
\setlist[enumerate,2]{labelindent=0pt}

% Blackboard Bold
\newcommand{\Z}{\mathbb{Z}}                 % integers
\newcommand{\Zpl}{\mathbb{Z}_{+}}           % positive integers
\newcommand{\Znn}{\mathbb{Z}_{\geq 0}}      % nonnegative integers. 
\newcommand{\Q}{\mathbb{Q}}                 % rationals
\newcommand{\Qpl}{\mathbb{Q}_{+}}           % positive Rationals
\newcommand{\R}{\mathbb{R}}                 % reals
\newcommand{\Rpl}{\mathbb{R}_{+}}           % positive Reals
\newcommand{\C}{\mathbb{C}}                 % complex numbers
\newcommand{\F}{\mathbb{F}}                 % field

% Words
\newcommand{\st}{\text{ such that }}
\newcommand{\wrt}{\text{ with respect to }}
\newcommand{\with}{\text{ with }}
\newcommand{\ie}{\text{, i.e. }}
\newcommand*{\ditto}{\text{---}\texttt{"}\text{---}}

% Operators
\newcommand{\abs}[1]{\left|#1\right|}                   % absolute value
\newcommand{\norm}[1]{\abs{\abs{#1}}}
\newcommand{\floor}[1]{\left\lfloor #1 \right\rfloor}   % floor
\newcommand{\ceil}[1]{\left\lceil #1 \right\rceil}      % ceiling
\newcommand{\powset}[1]{\mathscr{P}(#1)}

% Category Theory 
\renewcommand{\hom}{\text{Hom}}
\newcommand{\homspace}[3]{\hom_{#1}(#2, #3)}
\newcommand{\coproduct}[9]{
  \begin{tikzcd}[ampersand replacement=\&]
      \& \& #1 \arrow[dll, bend right, "#6"] \arrow[dl, "#5"] \\
   #4 \& #3 \arrow[l, "#9"] \\
      \& \& #2 \arrow[ull, bend left, "#8"] \arrow[ul, "#7"]
  \end{tikzcd}
}

\newcommand{\product}[9]{ 
  \begin{tikzcd}[ampersand replacement=\&]
      \& \& #3 \\
      #1 \arrow[r, "#7"] \arrow[urr, bend left, "#5"] \arrow[drr, bend right, "#6"] \& #2 \arrow[ur, "#8"] \arrow[dr, "#9"] \\ 
      \& \& #4
  \end{tikzcd}
}


% Algebra 
\newcommand{\Syl}[2]{\text{Syl}_{#1}{(#2)}}           % set of Sylow-p subgroups 
\newcommand{\<}{\langle}                % \<x\>, subgroup generated by x 
\renewcommand{\>}{\rangle}
\renewcommand{\index}[2]{[#1 : #2]}
\newcommand{\id}{\text{id}}             % identity element
\newcommand{\lhdneq}{%
  \mathrel{\ooalign{$\lneq$\cr\raise.22ex\hbox{$\lhd$}\cr}}}
\newcommand{\order}[1]{\text{o}(#1)}    % order of an element
\newcommand{\spec}{\operatorname{Spec}}
\newcommand{\rad}{\operatorname{rad}}   % radical of an ideal 
\newcommand{\sym}[1]{\mathfrak{S}_{#1}}
\newcommand{\aut}[1]{\text{Aut}(#1)} 
\newcommand{\gal}[2]{\text{Gal}(#1 / #2)} 
\newcommand{\inn}[1]{\text{Inn}(#1)} 
\let\oldcong\cong
\let\oldequiv\equiv
\renewcommand{\cong}{\oldequiv}
\renewcommand{\equiv}{\oldcong}

\newcommand{\triangleleftneq}{\mathrel{\ooalign{%
  $\trianglelefteq$\cr
  \hidewidth\raise0.06ex\hbox{$\not\mkern4mu$}\hidewidth\cr
}}}

\makeatletter % cycle
\newcommand{\cyc}[1]{(\mathbf{\cyc@process#1\relax})}
\def\cyc@process#1#2\relax{%
	#1%
	\ifx\relax#2\relax
	\else
		\,\cyc@process#2\relax
	\fi
}
\makeatother

\newcommand{\GL}[2]{\text{GL}_{#1}(#2)}                             % general linear group
\newcommand{\SL}[2]{\text{SL}_{#1}(#2)}                             % special linear group
\newcommand{\gl}[2]{\mathfrak{gl}_{#1}(#2)}
\renewcommand{\sl}[2]{\mathfrak{sl}_{#1}(#2)}
\newcommand{\so}[2]{\mathfrak{so}_{#1}(#2)}
\newcommand{\su}[1]{\mathfrak{su}(#1)}

% Linear Algebra
\newcommand{\bmat}[1]{\begin{bmatrix}#1\end{bmatrix}}   % bracketed matrix
\newcommand{\rank}{\operatorname{rank}}                 % rank
\newcommand{\nullity}{\operatorname{nullity}}           % nullity
\newcommand{\tr}{\operatorname{tr}}

% Combinatorics
\newcommand{\qnum}[1]{\left[#1\right]_q}                % q-analog
\newcommand{\qfact}[1]{\left[#1\right]!_q}              % q-analog factorial 
\newcommand{\qbinom}[2]{\binom{#1}{#2}_q}               % Gaussian binomial coefficient

% Topology/Analysis
\newcommand{\ball}[2]{\text{B}_{#1}(#2)}  % B_r(x): open r-balls around x
\newcommand{\diam}{\text{diam}}           % diamter of a set in metric space 
\newcommand{\lowint}[2]{\underline{\int_{#1}^{#2}}}
\newcommand{\upint}[2]{\overline{\int}_{#1}^{#2}}

% Misc Notation
\newcommand{\oneton}[1]{\{1, \ldots, #1\}}
\newcommand{\defeq}{\vcentcolon=}   % :=
\newcommand{\eqdef}{=\vcentcolon}   % =: 
\renewcommand{\bf}[1]{\textbf{#1}}
\newcommand{\im}{\operatorname{im}}
\newcommand{\fnrest}[2]{\left.#1\right|_{#2}}
\newcommand{\isom}{\overset{\sim}{\rightarrow}} 


\renewcommand{\thesection}{\arabic{section}.}
\renewcommand{\thesubsection}{(\alph{subsection})}

\newtheorem{claim}{Claim}
\newtheorem*{lemma}{Lemma}

%% Headers & title setup
\newcommand{\course}{Math100B}
\newcommand{\myname}{Jay Ser}
\setlength{\headheight}{14.5pt}
\pagestyle{fancy}
\fancyhf{}
\renewcommand{\headrulewidth}{0.4pt}
\lhead{\course}
\rhead{\myname}
\cfoot{\thepage}
\setlength{\droptitle}{-4em} 
\title{\course\ - HW \#8} 
\author{\myname}
\date{2026.03.08}

\begin{document}
\maketitle
\thispagestyle{fancy}

%------------------------------------------------------------------------------%

\section{}
Let $T_1(v_j) = \sum_{i = 1}^{n} a_{ij} v_j$ and $T_2(v_j) = \sum_{i = 1}^{n} b_{ij} v_j$. 
Then  
\begin{align*}
  (T_1 \circ T_2)(v_j) & = T_1 (\sum_{k = 1}^n b_{kj} v_k) \\ 
                       & = \sum_{k = 1}^{n} b_{kj} T_1 (v_k) \\ 
                       & = \sum_{k = 1}^{n} b_{kj} \sum_{i = 1}^{n} a_{ik}v_i \\ 
                       & = \sum_{i = 1}^{n} \left( \sum_{k = 1}^{n} a_{ik} b_{kj} \right)v_i 
\end{align*}
So the $i, j$th entry of $A_{T_1 \circ T_2, \mathcal{B}}$ is $\sum_{k = 1}^{n} a_{ik} b_{kj}$. 
This is exactly the $i, j$th entry of $A_{T_1, \mathcal{B}} A_{T_2, \mathcal{B}}$. 

\section{}
\subsection{}
Represent each $v \in V$ as a coordinate vector $\bmat{a_1, \ldots, a_n}^{T}$ where $v = \sum a_i v_i$. 
It suffices to check that $A_{S^{-1}, \mathcal{B}} A_{S, \mathcal{B}} (v_i) = v_i$ for all $1 \leq j \leq n$. 
For a given $j$, let $S(v_j) = w_j \defeq \sum_{i = 1}^{n} b_{ij} v_i$.
Then since $S^{-1}(w_j) = v_j$, $A_{S^{-1}, \mathcal{B}} (\bmat{b_{1j}, \ldots, b_{nj}}^T) = e_j$, the standard basis vector and the vector representation of $v_j$. 
Hence 
$$A_{S^{-1}, \mathcal{B}} A_{S, \mathcal{B}} (e_j) = A_{S^{-1}, \mathcal{B}} (\bmat{b_{1j}, \ldots, b_{nj}}^{T}) = e_j$$
In other words, $A_{S^{-1}, \mathcal{B}} A_{S, \mathcal{B}}$ is the identity matrix. 
This shows that $A_{S, \mathcal{B}}$ is invertible with inverse $A_{S^{-1}, \mathcal{B}}$. 

\subsection{} 
\newcommand{\PiAP}{A_{S, \mathcal{B}}^{-1} A_{T, \mathcal{B}} A_{S, \mathcal{B}}}
\newcommand{\Ap}{A_{T, \mathcal{B'}}} 
By Q1 and Q2 part (a), 
$$ \PiAP = A_{S^{-1} \circ T \circ S, \mathcal{B}}$$ 
For any $1 \leq j \leq n$, suppose $T(w_j) = \sum b_{ij} w_i$. 
Then, following the hint, 
$$S^{-1} \circ T \circ S(v_j) = \sum_{i = 1}^{n} b_{ij} v_i$$   
So the $i, j$th entry of $\PiAP$ is $b_{ij}$. 
By definition, the $i, j$th entry of $\Ap$ is $b_{ij}$. 
Hence 
$$\PiAP = \Ap$$ 

\section{} 
Let $\mathcal{B} = \{v_1, \ldots, v_n\}$ and $\mathcal{B'} = \{w_1, \ldots, w_n\}$ be two bases for $V$. 
Let $A$ and $A'$ be the matrix representation of the linear transformation $T$ with respect to $\mathcal{B}$ and $\mathcal{B'}$, respectively, and let $P$ be the transition matrix from $\mathcal{B}$ to $\mathcal{B'}$. 
By Q2 (b), $A' = P^{-1}AP$. 
Since $\det{AB} = \det{A} \det{B}$ for any matrices $A$ and $B$ over any field, 
$$\det{A'} = \det{P^{-1}} \det{A} \det{P} = \det{A}$$
where $\det{P}\det{P^{-1}} = 1$ since $P P^{-1} = I$. 
This shows that the determinant of a linear transformation is independent of the choice of basis. 

\section{} 
Fix $\alpha \in R$. 
Assume $\alpha$ is a unit in $R$.
Let $\alpha' = \alpha^{-1}$. 
Then $\varphi_{\alpha'} = \varphi_{\alpha}^{-1}$. 
Hence $\det \varphi_{\alpha} \det \varphi_{\alpha'} = 1$.
Equivalently, 
$$N(\alpha) N(\alpha') = 1$$
Hence $N(\alpha)$ is a unit in $F$. 
Conversely, if $\alpha$ is not a unit in $R$, then $\varphi_{\alpha}$ is not bijective because it is, in particular, not surjective: 
$$\varphi_{\alpha}(R) = \alpha R \subsetneq R$$
where $\alpha R$ is, by definition, the ideal generated by $\alpha$ in $R$. 
This means any matrix representation of $\varphi_{\alpha}$ will have linearly dependent rows/columns, hence $\det \varphi_\alpha = 0$;
$N(\alpha)$ is not a unit in $F$. 

\section{} 
\subsection{}
By definition, $\mathcal{B} = \{1, i\}$ spans $\Q[i]$. 
It is also clear that $\alpha 1 + \beta i = 0$ for $\alpha, \beta \in \Q \iff \alpha = \beta = 0$, so $\{1, i\}$ is a linearly independent set.  
This demonstrates that $\{1, i\}$ is a basis for $\Q[i]$. 

\subsection{}
$T(1) = 1(1) + 0(i)$ and $T(i) = 0(1) + (-1)i$, so 
$$A_{T, \mathcal{B}} = \bmat{1 & 0 \\ 0 & -1}, \: \det T = -1$$ 

\subsection{} 
Let $\varphi: R \rightarrow R; \: x \mapsto (a + bi)x$. 
Then $\varphi(1) = a + bi$ and $\varphi(i) = -b + ai$, so the matrix representation of $\varphi \wrt \mathcal{B}$ is $\bmat{a & -b \\ b & a}$, hence
$$N(\alpha) = \det \varphi = a^2 + b^2$$

\section{}
Let $\dim V_1 = n$, $\dim \ker T = k$.
Take $\mathcal{B} = \{v_1, \ldots, v_k, w_{k + 1}, \ldots, w_n\}$ such that $\mathcal{B}$ is a basis for $V_1$ and $\{v_1, \ldots, v_k\}$ is a basis for $\ker T$. 
Then $\{T(w_{k + 1}), \ldots, T(w_{n})\}$ spans $T(V_1)$ since given any $T(v) \in T(V_1)$, 
$$T(v) = \sum_{i = 1}^{k} a_i T(v_i) + \sum_{i = k + 1}^{n} a_i T(w_i) = \sum_{i = k + 1}^{n} a_i T(w_i)$$
where $v = \sum_{i = 1}^{k} a_i v_i + \sum_{i = k + 1}^{n} a_i w_i$. 
So the isomorphism $V_1 \equiv \ker T \times T(V_1)$ can be established by the linear transformation  
$$\varphi: V_1 \rightarrow \ker T \times T(V_1); \: \sum_{i = 1}^{k} a_i v_i + \sum_{i = k + 1}^{n} a_i w_i \mapsto \left(\sum_{i = 1}^{k} a_i v_i, \sum_{i = k + 1}^{n} a_i T(w_i) \right)$$
$\varphi$ is bijective by the analysis above. 
I hope it is not controversial to say without proof that $\varphi$ is indeed a linear transformation...

\section{}
\subsection{}
Yes. 

\subsection{}
By definition, 
$$T: 1 \mapsto 0, \: T: x^k \mapsto x^{k - 1} \text{ for } k \geq 1$$
So the matrix representation of $T \wrt \mathcal{B} = \{1, x, \ldots, x^d\}$ is 
$$A_{T, \mathcal{B}} = \bmat{0 & 1 \\ 0 & 0 & 1 \\ \vdots & & & \ddots \\ 0 & & \cdots & 0 & 1 & \\ 0 & & \cdots & & 0}$$ 
(I don't think I typeset this matrix very well. Sorry, Gavin!). 
$\det T = 0$ since $A_{T, \mathcal{B}}$ has a column of zeros. 

\subsection{}
Take any $f(x), g(x) \in \R[x]$ of degree at most $d$. 
Write $f(x) = a_d x^d + a_{d - 1}x^{d - 1} + \ldots + a_0$ and represent $f$ by the column vector $\alpha = \bmat{a_0 & a_1 & \cdots & a_d}^{T}$. 
Similarly, let the vector representation of $g$ be $\beta = \bmat{b_0, b_1, \ldots, b_d}^{T}$. 
Then 
$$A_{T, \mathcal{B}} \alpha = A_{T, \mathcal{B}} \beta \iff a_1 = b_1, \, a_2 = b_2, \, \ldots, a_n = b_n$$ 
In other words, a polynomial is uniquely determined by its derivative up to the constant term. 

\section{}
Let $f(x) = x^2 + bx + c$ and $g(x) = x^2 -(b^2 - 4c)$. 
Both polynomials are irreducible if and only if they have a root. 
Suppose $\alpha \in F$ is a root of $g$: 
$\alpha^2 = b^2 - 4c$.
Then 
\begin{align*}
  f(2^{-1}(-b + \alpha)) & = 4^{-1} (\alpha^2 - 2\alpha b + b^2) - 2^{-1} b^2 + 2^{-1} \alpha b + c \\ 
                         & = 4^{-1} (b^2 - 4c - 2\alpha b + b^2) - 2^{-1} b^2 + 2^{-1} \alpha b + c \\ 
                         & = 0
\end{align*}
so $f$ is reducible over $F$.  

Next, suppose $f$ is reducible with root $x$: 
$x^2 + bx + c = 0$. 
Then 
\begin{align*}
  g(2x + b) & = (2x + b)^2 - b^2 - 4c \\ 
            & = 4x^2 + 4xb + b^2 - b^2 - 4c \\ 
            & = 4(-bx - c) + 4xb -4c \\ 
            & = 0
\end{align*}
so $g$ is reducible over $F$ 

\end{document}
